\chapter{Statistiques : position et lecture critique}
\section*{Exercice 18.1 -- Indicateurs}
La somme vaut $52$, donc la moyenne $5{,}2$. La médiane est $(5+5)/2=5$. Les rangs des quartiles sont $\lceil10/4\rceil=3$ et $\lceil30/4\rceil=8$, donc $Q_1=2$, $Q_3=8$. L'étendue vaut $9$, l'écart interquartile $6$.

\section*{Exercice 18.2 -- Moyenne de groupes}
Les sommes sont $12\times14=168$ et $18\times11=198$. Leur total est $366$ pour $30$ élèves, donc la moyenne vaut $12{,}2$.

\section*{Exercice 18.3 -- Transformation}
La linéarité donne $\overline y=3\times7-2=19$. On peut le retrouver en sommant $3x_i-2$ : la somme devient trois fois l'ancienne somme moins deux fois l'effectif.

\section*{Exercice 18.4 -- Ajout extrême}
La somme devient $152$, l'effectif $11$, donc la moyenne vaut $152/11\simeq13{,}82$. La valeur centrale est de rang $6$, donc la médiane reste $5$. La valeur extrême déplace beaucoup plus la moyenne.

\section*{Exercice 18.5 -- Lire un diagramme}
La fréquence vaut $45/180=1/4$. Les angles sont $360/4=90\degre$ et $180/4=45\degre$. La relation vient de la proportionnalité de l'angle à l'aire du secteur pour un rayon fixé.

\section*{Exercice 18.6 -- Mêmes centres ?}
La moyenne et la médiane valent $5$ dans les deux séries, mais les étendues sont $2$ et $10$. Les distributions ne sont pas identiques : $B$ est plus dispersée. Deux indicateurs centraux ne déterminent pas une série.
